\begingroup
\setlength{\parindent}{0pt}
\setlength{\parskip}{4pt}
\colorlet{panelAccent}{fourierAccent}
\clearpage
\subsection{Input}

Patch a signal to an input to see its spectrum. The four inputs run from
red at the top through green, blue, and yellow. Each has its own trace and
gain knob, so you can compare up to four signals with the same analysis
settings. Leave the gains at \textbf{0 dB} for a fair level comparison.

Gain changes what Fourier analyzes; it does not change what you hear through
your existing audio path. Turning it fully down silences that input's
analysis. The maximum, +12 dB, is about four times the input amplitude.

A polyphonic cable is combined into one signal at each jack. Voices are not
shown separately: they can reinforce or cancel one another. Split voices
into separate cables if you want to compare them individually.

\subsection{Run}

\textbf{Lit: capture new audio. Dark: inspect the audio already captured.}
Press Run to stop following a changing sound, then experiment with the
analysis settings. Press it again to resume. Run starts enabled.

Stopping capture does not stop Fourier's analysis. The last unfinished
analysis may complete, and Average can continue to settle. This makes Run
useful for comparing windows on the same retained sound.

\begin{center}
\renewcommand{\arraystretch}{1.35}
\begin{tabularx}{\textwidth}{@{}p{.35\textwidth}X@{}}
\toprule
\textbf{While Run is off} & \textbf{What happens} \\
\midrule
Window, Average, Smooth & Reanalyze the retained audio. \\
Scales, bounds, Slope & Change how the result is displayed. \\
Gain, AC coupling & Resume capture to see their effect in the analysis. \\
Length & Clears the retained audio; resume to fill it again. \\
\bottomrule
\end{tabularx}
\end{center}

\paragraph{Try it}
Capture a steady oscillator, stop Run, and switch Window between Hann and
Flattop. Watch the width of the peak change even though the captured sound
has not changed. The next page explains why.

\clearpage
\subsection{Window Function}
A pure tone does not always appear as one thin line. The analyzer works on
short pieces of audio, and the edges of each piece can spread the tone into
nearby frequencies. This spreading is called \textit{leakage}. A window
softens the edges, trading a broader main peak for less energy around it.

\textbf{Start with Flattop} when comparing the heights of steady tones.
Try \textbf{Hann} when you want to distinguish nearby harmonics. Boxcar gives
a narrow main peak, but its leakage can hide quieter tones beside a loud one.
No window is best for every sound. The FFT measures at evenly spaced
frequencies called \textit{bins}; the plots below show what it sees.

\begin{figure}[!htp]
\centering
\input{WindowTradeoffs.tex}
\caption{Calculated responses to one tone halfway between FFT bins.
The dashed line marks the tone; dots mark measured bins. All three plots use
the same 2048-sample length and level scale. The ripples are leakage, not
additional tones.}
\end{figure}

\paragraph{Read the shape, not just the height}
The small ripples around a large peak can come from the analysis itself;
they are not necessarily extra notes or noise in the source. Compare windows
before drawing that conclusion. The module compensates for the window's
level reduction, but peak shape and the displayed noise level still change.

\paragraph{Try it}
With Average and Smooth off, move an oscillator's pitch slowly. Flattop
usually gives a steadier peak height as the tone passes between bins;
Hann gives a narrower peak. Fourier can also compare these windows while
Run is off. The complete window menu is on page~\pageref{sec:control-values}.

\clearpage
\subsection{Length}

\textbf{Length decides how much audio goes into one spectrum.} A longer
piece helps separate nearby pitches, especially in the bass. A shorter piece
follows a changing sound more closely. Start at the default \textbf{2048}
samples, then increase Length when two nearby tones merge into one peak.

\begin{figure}[!htp]
\centering
% Derived geometry at 48 kHz: frame duration N/fs and bin spacing fs/N.
\begin{tikzpicture}[x=1cm,y=1cm,font=\footnotesize]
  \node[anchor=west,font=\small\bfseries] at (0,3.65) {Length};
  \node[anchor=west,font=\small\bfseries] at (2,3.65) {Time in one frame};
  \node[anchor=west,font=\small\bfseries] at (8.3,3.65) {Frequency bins};
  \foreach \n/\intervals/\y in {1024/1/2.8,4096/4/1.4,16384/16/0} {
    \node[anchor=west] at (0,\y) {$N=\n$};
    \pgfmathsetmacro{\duration}{5.4*\intervals/16}
    \draw[fill=black!12,draw=black!60] (2,\y-.17) rectangle +(\duration,.34);
    \draw[black!40] (8.3,\y) -- (12.3,\y);
    \foreach \k in {0,...,\intervals} {
      \pgfmathsetmacro{\x}{8.3+4*\k/\intervals}
      \draw[black!70] (\x,\y-.17) -- (\x,\y+.17);
    }
  }
  \node[anchor=west] at (2,2.25) {$21.3\,\mathrm{ms}$};
  \node[anchor=west] at (2,.85) {$85.3\,\mathrm{ms}$};
  \node[anchor=west] at (2,-.55) {$341.3\,\mathrm{ms}$};
  \node[anchor=west] at (8.3,2.25) {$46.9\,\mathrm{Hz}$ apart};
  \node[anchor=west] at (8.3,.85) {$11.7\,\mathrm{Hz}$ apart};
  \node[anchor=west] at (8.3,-.55) {$2.93\,\mathrm{Hz}$ apart};
\end{tikzpicture}

\caption{The tradeoff at a Rack sample rate of 48 kHz. A longer frame includes
more time and gives more closely spaced frequency bins. Each row shows the
same time and frequency spans. Window choice also affects whether two tones
can be separated.}
\end{figure}

A \textit{bin} is one of the evenly spaced frequencies measured by the FFT.
More bins can reveal finer detail; enlarging the view alone cannot. Longer
frames also require more computation. Changing Length clears the captured
audio, so leave Run on long enough to collect a new frame.

\subsection{Hop}

\textbf{Hop decides how often to start another spectrum.} Its default is
\textbf{30 ms}. Lower values follow changes with more frequent analyses;
higher values reduce the analysis workload. The available range is 5--300 ms.

Length and Hop do different jobs. Imagine taking overlapping photographs:
Length is the time covered by each exposure, and Hop is the time between
exposures. Reducing Hop does not shorten the audio inside each frame.
If Hop is longer than a frame, gaps appear between the analyzed pieces and
a short event can be missed. Rack's display refresh also affects when you
see an update; Hop is not a promise of screen latency.

\clearpage
\subsection{Freq Scale}

Choose \textbf{Linear} to compare harmonic spacing: equal steps in hertz
occupy equal widths. \textbf{Logarithmic} (the default) gives the bass more
room. This version uses a square-root mapping, so octaves are not equally
wide; read the labels or hover readout for exact frequencies.

\subsection{Y Scale}

Start with \textbf{Log 60dB} for everyday use. Choose \textbf{Log 120dB}
to see quieter detail. Both extend up to +12 dB; their names identify the
lower limit. \textbf{Linear} emphasizes large amplitudes, expressed as
percentages. Changing scale changes the view, not the sound.

\subsection{Average}

Increase \textbf{Average} to calm a flickering trace and compare broad trends.
Leave it at \textbf{0 ms} to follow attacks. Higher settings give previous
results more influence; they do not select a fixed-length recording. The
range extends to 2500 ms.

\subsection{Smooth}

Leave \textbf{Smooth} at \textbf{None} for individual harmonics. Try
\textbf{1/12 octave} for a gentler outline of tonal balance. Wider settings
merge neighboring peaks and can lower them; they do not add resolution.
The octave-based width covers more hertz in the treble than in the bass.

\begin{figure}[!htp]
\centering
\input{Smoothing.tex}
\caption{Average softens changes across time; Smooth blends detail across
frequency. These are qualitative sketches, not measured responses.}
\end{figure}

\clearpage
\subsection{Low Frequency}

\textbf{LO Freq} sets the left edge of the plot. Raise it to leave frequencies
below your area of interest out of the view. It starts at \textbf{0 Hz}.
This is a viewing control, not a high-pass filter on the sound.

\subsection{High Frequency}

\textbf{HI Freq} sets the right edge. Lower it to give the bass or a group of
harmonics more room. Keep it above LO Freq. Both controls can reach half of
Rack's sample rate, called the \textit{Nyquist frequency}; that is 24 kHz
at a 48 kHz sample rate. HI Freq starts at this upper limit.

\paragraph{Try it}
Set LO Freq to 0 Hz and HI Freq to 500 Hz to inspect the bass. The existing
peaks spread across a wider area, but no new frequency detail is added.
If nearby tones still merge, increase Length. Check your bounds after
changing Rack's sample rate, because the available range changes too.

\subsection{Slope}

\textbf{Slope tilts the view, like a display-only tone control.} Positive
values lift high frequencies and reduce low frequencies around a 1 kHz
pivot. Negative values reverse that tilt. Your audio is unaffected.

The default is \textbf{+4.5 dB per octave}, which can help show high-frequency
content. Set it to \textbf{0} when comparing unweighted levels. Its range is
$-9$ to $+9$ dB per octave.

\begin{center}
\renewcommand{\arraystretch}{1.3}
\begin{tabular}{@{}lrrr@{}}
\toprule
\textbf{Example: +3 dB/oct} & \textbf{500 Hz} & \textbf{1 kHz} & \textbf{2 kHz} \\
\midrule
Change in displayed level & $-3$ dB & 0 dB & $+3$ dB \\
\bottomrule
\end{tabular}
\end{center}

An octave doubles the frequency. The same tilt continues above and below
the example: it is not an equalizer acting on the source. Input gain,
windowing, and smoothing can also change the displayed levels; keep those
settings fixed when judging what Slope does.

\endgroup
